\section{Сжатие контекста не равно экономии на задаче}
\label{sec:economics}
Состояние модели сокращает контекст, но не сохраняет результат автоматически. На общих успехах медианное отношение пикового промпта к контролю равно \ModelStatePeakRatio{}, стоимости --- \ModelStateCostRatio{}, выходных токенов --- \ModelStateOutputRatio{}. По всем включённым исходам успешность снижается с \ModelControlRate\% до \ModelStateRate\%; бутстрэп-интервал разности равен $\ModelStateDeltaCI$ процентного пункта (дополнение~\ref{app:contrasts}).

Cache hit обнаруживает ещё одну цену замены истории. Медианная доля входа, обслуженного из кеша, составляет \ModelControlCacheHit\% для полной истории, \ModelStateCacheHit\% для состояния без истории, \ModelHistoryTwoCacheHit\% при $k=2$ и \ModelHistorySixCacheHit\% при $k=6$. Растущий транскрипт может повторно использовать неизменный префикс; обновляемое состояние и движущееся окно могут сокращать такое использование. Измерения подтверждают различие поведения кеша, но не выделяют конкретную деталь сериализации как его причину.

Сам по себе низкий cache hit не делает политику невыгодной. Окно с наименьшим попаданием в кеш одновременно даёт наибольшую успешность среди проверенных вариантов состояния. Полная цена складывается из кешированного и некешированного входа, генерации и числа действий. Для состояния без истории стоимость на наблюдаемый успех равна \$\ModelStateCostPerSuccess{} против \$\ModelControlCostPerSuccess{} у контроля. В дальнейших сравнениях сохраняются одинаковые колонки качества, денег, токенов и кеша: изменение представления памяти нужно оценивать по всем этим измерениям.
